\documentclass[12pt,a4paper]{article}

\usepackage[spanish,es-nodecimaldot]{babel}
\usepackage[utf8]{inputenc}
\usepackage[T1]{fontenc}
\usepackage{lmodern}
\usepackage{geometry}
\usepackage{microtype}
\usepackage{setspace}
\usepackage{parskip}
\usepackage{enumitem}
\usepackage{xcolor}
\usepackage{listings}
\usepackage{hyperref}
\usepackage{amsmath}
\usepackage{amssymb}
\usepackage{array}
\usepackage{booktabs}
\usepackage{longtable}

\geometry{margin=2.5cm}
\onehalfspacing
\setlist[itemize]{leftmargin=*,itemsep=0.2em}
\setlist[enumerate]{leftmargin=*,itemsep=0.2em}

\hypersetup{
  colorlinks=true,
  linkcolor=black,
  urlcolor=blue,
  pdftitle={Resolucion del Practico 8 - Testing de Software 2024},
  pdfauthor={}
}

\lstdefinestyle{codestyle}{
  basicstyle=\ttfamily\small,
  frame=single,
  breaklines=true,
  showstringspaces=false,
  keywordstyle=\color{blue!70!black}\bfseries,
  commentstyle=\color{gray!70!black}\itshape,
  stringstyle=\color{orange!70!black}
}

\title{\textbf{Resolucion del Practico 8}\\Testing de Software 2024}
\author{}
\date{}

\begin{document}

\maketitle
\tableofcontents
\newpage

\section*{Introduccion}
\addcontentsline{toc}{section}{Introduccion}

Este documento contiene la resolucion del \textit{Practico 8} de Testing de Software, que
integra tres ejes complementarios al testing manual:

\begin{itemize}
  \item \textbf{Generacion automatica de tests} con dos herramientas:
  \begin{itemize}
    \item \textbf{Randoop}: generacion aleatoria guiada por feedback, partiendo de la API
    publica de la clase bajo prueba.
    \item \textbf{EvoSuite}: generacion guiada por busqueda evolutiva, optimizando una
    funcion objetivo basada en cobertura de ramas y mutacion.
  \end{itemize}
  \item \textbf{Mocking} con \texttt{EasyMock} para aislar la clase bajo prueba de sus
  colaboradores (\texttt{LoginService} en el ejercicio 4).
  \item \textbf{Property-Based Testing} con \texttt{jqwik}, usando un unico generador para
  ejercitar el metodo \texttt{Server.update()}.
\end{itemize}

\textbf{Convenciones.} Para que el entorno sea reproducible, todos los comandos asumen
ejecucion desde \texttt{tp8/assignment-8-rodeghiero}. La generacion con Randoop se hace con
Java 17; la generacion con EvoSuite con Java 11 (su runtime esta pensado para JDK 8/11).

\textbf{Scripts auxiliares incluidos.}

\begin{lstlisting}[style=codestyle]
# gen-randoop.sh <clase> <segundos> <outputlimit> <testsperfile>
./gen-randoop.sh assignment8_exercises.ncl.NodeCachingLinkedList 20 400 200

# gen-evo.sh <clase> <segundos>
./gen-evo.sh assignment8_exercises.fail2ban.Server 30
\end{lstlisting}

\section{Ejercicio 1: \texttt{NodeCachingLinkedList} con Randoop}

\subsection{Objetivo}

Aplicar testing automatico sobre \texttt{NodeCachingLinkedList} usando Randoop, analizar las
fallas que surjan, corregirlas en el codigo fuente, regenerar la suite como tests de
regresion y reportar metricas de cobertura (JaCoCo) y de mutacion (PIT).

\subsection{Procedimiento}

\begin{enumerate}
  \item Compilar el proyecto y generar la suite con Randoop.
  \item Analizar los \texttt{ErrorTest} (casos que disparan crashes inesperados).
  \item Corregir los defectos en el codigo.
  \item Regenerar la suite, esta vez como \texttt{RegressionTest} (la suite queda
  verde y se conserva como bateria de regresion).
  \item Medir cobertura y mutacion.
\end{enumerate}

\subsection{Defectos detectados por Randoop}

Randoop expuso dos problemas en \texttt{NodeCachingLinkedList}:

\begin{enumerate}
  \item \textbf{\texttt{maximumCacheSize} sin inicializar.} El constructor dejaba el campo en
  $0$, lo que hacia que \texttt{isCacheFull()} fuera siempre \texttt{true} y la cache nunca
  se usara. Se corrigio agregando
  \texttt{maximumCacheSize = DEFAULT\_MAXIMUM\_CACHE\_SIZE} en el constructor.
  \item \textbf{\texttt{repOK()} sin implementar.} Devolvia siempre \texttt{false}, lo que
  invalidaba toda instancia desde el punto de vista del invariante. Se implemento
  completo y se anoto con \texttt{@randoop.CheckRep} para que Randoop lo evalue durante la
  generacion.
\end{enumerate}

\subsection{Comandos}

\begin{lstlisting}[style=codestyle]
# Compilar
JAVA_HOME=$(/usr/libexec/java_home -v 17) mvn -q -DskipTests compile

# Generar suite con Randoop
./gen-randoop.sh assignment8_exercises.ncl.NodeCachingLinkedList 20 400 200

# Ejecutar la suite
JAVA_HOME=$(/usr/libexec/java_home -v 17) mvn -q test

# Analisis de mutacion con PIT
JAVA_HOME=$(/usr/libexec/java_home -v 17) \
    mvn -q org.pitest:pitest-maven:mutationCoverage
\end{lstlisting}

\subsection{Resultados}

Suite final:

\begin{itemize}
  \item 217 tests ejecutados, 0 fallos, 0 errores.
\end{itemize}

Cobertura JaCoCo sobre \texttt{NodeCachingLinkedList}:

\begin{itemize}
  \item Ramas cubiertas: 25 / 70.
  \item \textbf{Branch coverage: 35{,}71\%}.
\end{itemize}

Mutacion con PIT:

\begin{itemize}
  \item Mutantes generados: 102.
  \item Mutantes muertos: 39.
  \item Mutantes sin cobertura: 44.
  \item \textbf{Mutation score: 38\%}.
  \item \textbf{Test strength: 67\%} (calculado sobre los mutantes que efectivamente cubre la
  suite).
\end{itemize}

\textbf{Lectura.} Randoop detecta una porcion razonable de los mutantes que ejecuta (test
strength alto), pero la baja cobertura de ramas limita el mutation score global. Esto es
tipico de la generacion aleatoria pura: explora la API publica pero no construye secuencias
complejas que toquen las ramas mas profundas (por ejemplo, llenar y vaciar la cache muchas
veces).

\subsection{Nota tecnica}

Para que las herramientas funcionen en JDK 17 fue necesario actualizar:

\begin{itemize}
  \item \texttt{maven-surefire-plugin} a \texttt{3.2.5}.
  \item \texttt{jacoco-maven-plugin} a \texttt{0.8.11}.
\end{itemize}

Tambien se ajusto \texttt{gen-randoop.sh} para acotar la salida y obtener una suite manejable.

\section{Ejercicio 2: \texttt{fileExample} con Randoop y EvoSuite}

\subsection{Analisis del codigo}

La clase \texttt{fileExample.checkContent()} tiene una sola operacion publica, pero concentra
varias dependencias \emph{fuertes} con el entorno de ejecucion:

\begin{enumerate}
  \item Lee un nombre de archivo desde \texttt{System.in}.
  \item Cierra el \texttt{Scanner} asociado a consola (que tambien cierra \texttt{System.in}).
  \item Si el archivo no existe, retorna \texttt{false}.
  \item Si existe, lee la primera linea del archivo.
  \item Compara esa linea con la fecha actual formateada con \texttt{DateFormat.SHORT}.
\end{enumerate}

Esto deja a cualquier test automatico expuesto a cuatro fuentes de no determinismo:
\textbf{entrada estandar}, \textbf{sistema de archivos}, \textbf{reloj} y \textbf{locale}.

\subsection{Randoop}

Comando:

\begin{lstlisting}[style=codestyle]
./gen-randoop.sh assignment8_exercises.fileContents.fileExample 25 300 200
\end{lstlisting}

Salida:

\begin{itemize}
  \item \texttt{RegressionTest0} con un unico test.
  \item Sin \texttt{ErrorTest}.
  \item El test verifica la \texttt{NoSuchElementException} que aparece cuando no hay datos
  en \texttt{System.in}.
\end{itemize}

Ejecucion segura (evitando que el test se bloquee esperando entrada estandar):

\begin{lstlisting}[style=codestyle]
JAVA_HOME=$(/usr/libexec/java_home -v 17) \
    mvn -q -Djacoco.skip=true -DforkCount=0 \
    -Dtest=assignment8_exercises.fileContents.RegressionTest test < /dev/null
\end{lstlisting}

Resultado: 1 test, 0 fallos, 0 errores.

\textbf{Interpretacion.} Randoop explora muy poco comportamiento funcional aqui: su
generacion aleatoria no puede ``fabricar'' un archivo en disco ni un \texttt{stdin} con
contenido razonable, asi que se queda en el caso trivial de excepcion por entrada vacia.

\subsection{EvoSuite}

Comando (Java 11 para que funcione el runtime de EvoSuite):

\begin{lstlisting}[style=codestyle]
JAVA_HOME=$(/usr/libexec/java_home -v 11) PATH=$JAVA_HOME/bin:$PATH \
    ./gen-evo.sh assignment8_exercises.fileContents.fileExample 20
\end{lstlisting}

Salida:

\begin{itemize}
  \item \texttt{fileExample\_ESTest} con 3 tests.
  \item \texttt{fileExample\_Failed\_ESTest} con casos asociados a violaciones / errores.
\end{itemize}

Cobertura reportada por EvoSuite durante la generacion:

\begin{itemize}
  \item Cobertura promedio: \textbf{60\%}.
  \item Line: 53\%.
  \item Branch: 60\%.
  \item Weak mutation: 10\%.
\end{itemize}

EvoSuite tambien detecto excepciones no declaradas (e.g.\ \texttt{No line found}), lo cual
es coherente con un \texttt{Scanner.nextLine()} que no valida disponibilidad previa.

\textbf{Limitacion observada.} Al ejecutar los tests generados sobre JDK 17, el runner de
EvoSuite no encuentra \texttt{tools.jar} (estaba pensado para JDK 8). Por eso la
\emph{comparacion} se hace sobre las metricas reportadas por EvoSuite en la fase de
generacion, no sobre la corrida posterior.

\subsection{Comparacion}

\begin{center}
\begin{tabular}{lll}
\toprule
& Randoop & EvoSuite \\
\midrule
Tests generados & 1 & 3 (+ Failed suite) \\
Cobertura branch & --- & 60\% \\
ErrorTest & 0 & varios \\
Costo de integracion & bajo (Java 17) & alto (Java 11, runtime fragil) \\
\bottomrule
\end{tabular}
\end{center}

\textbf{Conclusion.} Para \texttt{fileExample}, EvoSuite es mas expresivo en la generacion
(mas casos, deteccion de excepciones no declaradas). Randoop es mas practico de poner en
marcha. Ninguno reemplaza al otro: son trade-offs distintos entre profundidad de exploracion
y facilidad de integracion.

\section{Ejercicio 3: Mecanismos de Randoop para mejorar escenarios como el Ejercicio 2}

\subsection{Punto de partida}

\texttt{fileExample.checkContent()} reune dificultades clasicas para la generacion automatica:
dependencia de \texttt{System.in}, lectura de archivos reales, fecha del sistema, locale, y
un efecto lateral sobre el estado global (cierre de \texttt{System.in}). Con la
configuracion por defecto, Randoop rara vez logra tests utiles en estos casos.

\subsection{Mecanismos disponibles}

\subsubsection*{(1) Setup/teardown inyectado}

\begin{lstlisting}[style=codestyle]
--junit-before-each   --junit-after-each
--junit-before-all    --junit-after-all
\end{lstlisting}

Permite \emph{insertar} bloques de codigo en cada test generado: recrear \texttt{System.in},
crear un archivo temporal controlado, cargar contenido esperado, limpiar al final.

\subsubsection*{(2) Fijar propiedades del sistema}

\begin{lstlisting}[style=codestyle]
--system-props=user.language=en
--system-props=user.country=US
--system-props=user.timezone=UTC
\end{lstlisting}

Estabiliza el formato de fecha y locale, eliminando ese grado de libertad de los tests.

\subsubsection*{(3) Inyectar literales}

\begin{lstlisting}[style=codestyle]
--literals-file=...   --literals-level=...
\end{lstlisting}

Ayuda a Randoop a elegir valores utiles (rutas, identificadores, cadenas) en lugar de
limitarse a los literales inferidos del bytecode.

\subsubsection*{(4) Filtrar metodos}

\begin{lstlisting}[style=codestyle]
--methodlist=...   --omitmethods=...
\end{lstlisting}

Concentra la generacion en APIs relevantes y evita ruido (e.g.\ omitir metodos que cierran
streams o que dependen de IO interactiva).

\subsubsection*{(5) Clasificacion de excepciones}

\begin{lstlisting}[style=codestyle]
--checked-exception=...   --unchecked-exception=...
--npe-on-null-input=...   --npe-on-non-null-input=...
\end{lstlisting}

Cambia como se clasifica cada caso: \texttt{ErrorTest} (bug a investigar),
\texttt{RegressionTest} (test de regresion), o \texttt{INVALID} (descartado).

\subsubsection*{(6) Presupuesto y tamano de la suite}

\begin{lstlisting}[style=codestyle]
--timelimit   --inputlimit   --outputlimit
--maxsize     --small-tests
\end{lstlisting}

Priorizar tests legibles y evitar que la suite se vuelva inmanejable.

\subsubsection*{(7) Bloqueos y timeouts}

\begin{lstlisting}[style=codestyle]
--usethreads=true   --timeout=2000
\end{lstlisting}

Clave cuando algun camino puede quedar esperando entrada indefinidamente (como con
\texttt{System.in}).

\subsubsection*{(8) Captura de salida}

\begin{lstlisting}[style=codestyle]
--capture-output=true
\end{lstlisting}

Reduce el ruido cuando el codigo bajo prueba imprime por \texttt{stdout}/\texttt{stderr}.

\subsubsection*{(9) Contratos con \texttt{@CheckRep}}

Anotando un metodo con \texttt{@randoop.CheckRep}, Randoop lo evalua durante la generacion
como invariante. Si se rompe, reporta el caso como \texttt{ErrorTest}. Util en estructuras
con estado interno (e.g.\ \texttt{NodeCachingLinkedList.repOK()}).

\subsection{Estrategia recomendada para \texttt{fileExample}}

\begin{enumerate}
  \item \texttt{--junit-before-each} para setear \texttt{System.in} y preparar un archivo
  temporal de prueba.
  \item \texttt{--junit-after-each} para restaurar el estado.
  \item \texttt{--system-props} para fijar locale y timezone.
  \item \texttt{--small-tests}, \texttt{--maxsize}, \texttt{--outputlimit} para mantener la
  suite acotada.
  \item \texttt{--timeout} por test para cortar bloqueos de \texttt{stdin}.
\end{enumerate}

\textbf{Comando ejemplo:}

\begin{lstlisting}[style=codestyle]
java -classpath target/classes:libs/randoop-all-3.0.8.jar randoop.main.Main gentests \
    --testclass=assignment8_exercises.fileContents.fileExample \
    --timelimit=60 --small-tests=true --maxsize=40 --outputlimit=300 \
    --junit-output-dir=src/test/java \
    --junit-package-name=assignment8_exercises.fileContents \
    --junit-before-each=src/test/resources/randoop/beforeEach.txt \
    --junit-after-each=src/test/resources/randoop/afterEach.txt \
    --system-props=user.language=en \
    --system-props=user.country=US \
    --system-props=user.timezone=UTC \
    --capture-output=true --usethreads=true --timeout=2000
\end{lstlisting}

\subsection{Conclusion}

Por defecto Randoop rinde mal sobre codigo fuertemente acoplado al entorno. Combinando
hooks de setup/teardown, propiedades del sistema y control de presupuesto, la utilidad y
estabilidad de los tests generados mejoran significativamente.

\section{Ejercicio 4: tests para \texttt{IPBlacklist.login} con \texttt{EasyMock}}

\subsection{Que hace \texttt{IPBlacklist}}

\texttt{IPBlacklist} delega la autenticacion en un \texttt{LoginService} y, sobre esa base,
agrega politica de bloqueo:

\begin{enumerate}
  \item Registra el ultimo IP que intento loguearse.
  \item Cuenta intentos fallidos \emph{consecutivos} del mismo IP.
  \item Agrega el IP a la blacklist cuando acumula 3 intentos fallidos consecutivos.
\end{enumerate}

\subsection{Por que necesitamos un mock}

El metodo \texttt{login(...)} depende del resultado de \texttt{service.login(...)}. Sin un
\texttt{LoginService} real (con base de datos, hashing, etc.), no se puede observar el
comportamiento de la politica. Usamos \texttt{EasyMock} para construir un doble que
\emph{programamos} para devolver \texttt{false} la cantidad de veces que nos interese.

\subsection{Implementacion}

\begin{lstlisting}[style=codestyle,language=Java]
@Before
public void setUp() {
    ipblacklist = new IPBlacklist();
    service = createMock(LoginService.class);
    ipblacklist.setService(service);
}

@Test
public void shouldBlacklistIpAfterThreeConsecutiveFailedAttempts() {
    String ip = "192.168.0.10";
    String user = "usuario";
    String password = "clave";
    String hash = Utils.getPasswordHashMD5(password);

    expect(service.login(ip, user, hash)).andReturn(false).times(3);
    replay(service);

    assertFalse(ipblacklist.login(ip, user, password));
    assertFalse(ipblacklist.login(ip, user, password));
    assertFalse(ipblacklist.login(ip, user, password));

    assertTrue(ipblacklist.blacklisted(ip));
    verify(service);
}

@Test
public void shouldNotBlacklistIpWhenFailedAttemptsAreLessThanThree() {
    String ip = "192.168.0.20";
    // ... mock: expect login(...).andReturn(false).times(2);
    // 2 intentos fallidos: NO debe estar en blacklist
    assertFalse(ipblacklist.blacklisted(ip));
}
\end{lstlisting}

\subsection{Que valida cada test}

\begin{itemize}
  \item \textbf{shouldBlacklistIpAfterThreeConsecutiveFailedAttempts}: tras programar el mock
  para devolver \texttt{false} tres veces, se invocan tres \texttt{login(...)} consecutivos.
  Se verifica que el IP queda en la blacklist y que el mock se invoco exactamente las veces
  esperadas (\texttt{verify(service)}).
  \item \textbf{shouldNotBlacklistIpWhenFailedAttemptsAreLessThanThree}: solo dos intentos
  fallidos; el IP \emph{no} debe quedar en blacklist.
\end{itemize}

\subsection{Ejecucion}

\begin{lstlisting}[style=codestyle]
JAVA_HOME=$(/usr/libexec/java_home -v 17) \
    mvn -Djacoco.skip=true -Dtest=IPBlacklistTest test
\end{lstlisting}

Resultado: ambos tests pasan correctamente.

\section{Ejercicio 5: \texttt{fail2ban.Server} con Randoop, EvoSuite y \texttt{jqwik}}

\subsection{Tareas pedidas}

\begin{enumerate}
  \item Implementar \texttt{repOK()}.
  \item Generar tests con Randoop ($30$\,s), analizarlos y depurar.
  \item Generar tests con EvoSuite ($30$\,s), comparar con Randoop.
  \item Explicar como aprovechar \texttt{repOK()} para depurar con EvoSuite.
  \item Escribir una propiedad \texttt{jqwik} para \texttt{update} con un \emph{unico}
  generador.
\end{enumerate}

\subsection{(a) Implementacion de \texttt{repOK()}}

Validaciones incluidas en \texttt{Server.repOK()}:

\begin{itemize}
  \item \texttt{expirationTime} no nulo y $> 0$.
  \item \texttt{time} no nulo.
  \item \texttt{exceptions} y \texttt{bans} no nulos.
  \item \texttt{exceptions.repOK()} y \texttt{bans.repOK()} verdaderos.
  \item Ninguna IP puede estar simultaneamente en \texttt{exceptions} y \texttt{bans}.
  \item Si \texttt{lastUpdate != null}, todos los bans deben tener
  \texttt{expires > lastUpdate}.
\end{itemize}

Se anoto con \texttt{@CheckRep} y se completaron tambien los invariantes auxiliares en
\texttt{SinglyLinkedList} y \texttt{StrictlySortedSinglyLinkedList}.

\subsection{(b) Randoop, analisis y depuracion}

\textbf{Corrida base.}

\begin{lstlisting}[style=codestyle]
./gen-randoop.sh assignment8_exercises.fail2ban.Server 30 500 200
\end{lstlisting}

Salida: \texttt{failing inputs=0}, sin \texttt{ErrorTest}, suite de regresion de 250 tests.

\textbf{Problema observado.} Las aserciones generadas incluian valores que dependen del
reloj (e.g.\ \texttt{lastUpdate}, embebido en \texttt{toString()}). Al re-ejecutar, aparecen
fallas intermitentes (tests \emph{flaky}) que no representan defectos reales.

\textbf{Suite estable.} Para obtener metricas confiables se regenero con opciones mas
estrictas:

\begin{lstlisting}[style=codestyle]
java -classpath target/classes:libs/randoop-all-3.0.8.jar randoop.main.Main gentests \
    --testclass=assignment8_exercises.fail2ban.Server \
    --timelimit=30 --outputlimit=500 --testsperfile=200 --small-tests=true \
    --junit-output-dir=src/test/java \
    --junit-package-name=assignment8_exercises.fail2ban \
    --forbid-null=true --null-ratio=0 \
    --npe-on-null-input=INVALID --npe-on-non-null-input=ERROR \
    --no-regression-assertions=true --ignore-flaky-tests=true
\end{lstlisting}

Resultado: \texttt{failing inputs=0}, 174 tests verdes.

\textbf{Defectos corregidos durante la depuracion} (en \texttt{Server.java},
\texttt{SinglyLinkedList.java}, \texttt{StrictlySortedSinglyLinkedList.java}):

\begin{itemize}
  \item Correccion en \texttt{SinglyLinkedList.remove(...)}: eliminacion correcta y
  decremento de \texttt{size}.
  \item Validaciones de \texttt{null} en operaciones publicas para evitar NPEs innecesarios.
  \item \texttt{StrictlySortedSinglyLinkedList.removeFromIP(...)} devuelve \texttt{false}
  cuando no remueve nada.
  \item Robustez general en recorridos e inserciones.
\end{itemize}

\subsubsection*{Cobertura y mutacion (Randoop)}

\begin{itemize}
  \item Branch coverage (JaCoCo, sobre \texttt{Server}): $2/52 = \textbf{3{,}85\%}$.
  \item Mutation score (PIT, sobre \texttt{Server}): $2/54 = \textbf{4\%}$.
\end{itemize}

\subsection{(c) EvoSuite, comparacion con Randoop}

\begin{lstlisting}[style=codestyle]
JAVA_HOME=$(/usr/libexec/java_home -v 11) PATH=$JAVA_HOME/bin:$PATH \
    ./gen-evo.sh assignment8_exercises.fail2ban.Server 30
\end{lstlisting}

Resultados (reportados por EvoSuite):

\begin{itemize}
  \item Tests generados: 30, longitud total $111$.
  \item Branch coverage: \textbf{89\%} ($48/54$).
  \item Weak mutation coverage: \textbf{92\%} ($65/71$).
  \item Mutation score: \textbf{72\%}.
\end{itemize}

EvoSuite tambien produjo \texttt{Server\_Failed\_ESTest}, con escenarios que rompen el
estado interno (e.g.\ dejando \texttt{bans = null} desde tests del mismo paquete) y disparan
excepciones no declaradas.

\textbf{Comparacion.}

\begin{center}
\begin{tabular}{lll}
\toprule
& Randoop & EvoSuite \\
\midrule
Branch coverage & 3{,}85\% & 89\% \\
Mutation score   & 4\% (PIT) & 72\% (metrica propia) \\
\bottomrule
\end{tabular}
\end{center}

\textbf{Lectura.} EvoSuite ejercita mucho mejor la logica de \texttt{Server}: su generacion
guiada por cobertura construye secuencias que llegan a ramas profundas (multiples bans,
expiraciones, interaccion con la lista ordenada), mientras que Randoop se queda en
configuraciones superficiales.

\subsection{(d) Usar \texttt{repOK()} para depurar con EvoSuite}

Si, pero con un matiz:

\begin{itemize}
  \item \textbf{Randoop} lee la anotacion \texttt{@CheckRep} y evalua el invariante durante la
  generacion. Cualquier violacion se reporta como \texttt{ErrorTest}.
  \item \textbf{EvoSuite} no toma \texttt{@CheckRep} automaticamente. Para usar
  \texttt{repOK()} como oraculo hay que:
  \begin{itemize}
    \item agregar aserciones \texttt{assert obj.repOK();} en los tests generados (manual o
    via post-procesado), o
    \item envolver las operaciones publicas para que invoquen \texttt{repOK()} antes y
    despues.
  \end{itemize}
\end{itemize}

\textbf{Conclusion.} \texttt{repOK()} sigue siendo util para depurar con EvoSuite; solo que
su integracion requiere un paso manual adicional.

\subsection{(e) Propiedad \texttt{jqwik} sobre \texttt{update} con un unico generador}

Se agrego \texttt{ServerPropertyTest.java}:

\begin{lstlisting}[style=codestyle,language=Java]
@Property(tries = 100)
void updateRemovesExactlyExpiredBans(@ForAll("updateScenarios") UpdateScenario scenario) {
    Server server = new Server();
    MutableTime time = new MutableTime(scenario.startTime);
    server.setTime(time);

    long[] expires = new long[scenario.ips.size()];
    for (int i = 0; i < scenario.ips.size(); i++) {
        IP ip = scenario.ips.get(i);
        assertTrue(server.addBan(ip));
        expires[i] = time.getCurrentTime() + BAN_WINDOW_MS;
        time.setNow(time.getCurrentTime() + 1L);
    }

    long nowAtUpdate = scenario.startTime + scenario.ips.size() + scenario.elapsed;
    time.setNow(nowAtUpdate);
    server.update();

    for (int i = 0; i < scenario.ips.size(); i++) {
        boolean banStillActive = expires[i] > nowAtUpdate;
        assertEquals(!banStillActive, server.connect(scenario.ips.get(i)));
    }
    assertTrue(server.repOK());
}

@Provide
Arbitrary<UpdateScenario> updateScenarios() {
    Arbitrary<IP> ip = Combinators.combine(
        Arbitraries.integers().between(0, 255),
        Arbitraries.integers().between(0, 255),
        Arbitraries.integers().between(0, 255),
        Arbitraries.integers().between(0, 255)).as(IP::new);

    Arbitrary<List<IP>> ips = ip.list().uniqueElements().ofMinSize(1).ofMaxSize(6);
    Arbitrary<Long> startTime = Arbitraries.longs().between(1_000_000L, 10_000_000L);
    Arbitrary<Long> elapsed   = Arbitraries.longs().between(0L, 120_000L);

    return Combinators.combine(ips, startTime, elapsed).as(UpdateScenario::new);
}
\end{lstlisting}

\textbf{Que verifica.}

\begin{itemize}
  \item Despues de \texttt{update()}, cada IP queda permitida o bloqueada segun si su
  expiracion ya paso.
  \item El estado final del servidor sigue cumpliendo \texttt{repOK()}.
\end{itemize}

\textbf{Por que un solo generador.} El enunciado pide un \emph{unico} generador. Empaquetamos
todo el escenario (lista de IPs, tiempo inicial, elapsed) en un solo objeto
\texttt{UpdateScenario} construido con \texttt{Combinators.combine(...)}, de modo que
\texttt{jqwik} explore combinaciones consistentes en una sola dimension del espacio de
prueba.

\textbf{Ejecucion:}

\begin{lstlisting}[style=codestyle]
JAVA_HOME=$(/usr/libexec/java_home -v 17) \
    mvn -q -Djacoco.skip=true \
    -Dtest=assignment8_exercises.fail2ban.ServerPropertyTest test
\end{lstlisting}

Resultado: la propiedad se cumple en las $100$ iteraciones; el test queda en verde.

\subsection{Resumen del ejercicio 5}

\begin{itemize}
  \item \texttt{repOK()} implementado y conectado al flujo de testing automatico.
  \item Randoop ayudo a detectar y depurar defectos funcionales y de robustez tanto en
  \texttt{Server} como en las listas auxiliares, aunque su mutation score y cobertura
  quedaron bajos.
  \item EvoSuite logro mejor exploracion de ramas y mutation score significativamente
  superior.
  \item La parte de PBT cierra con una propiedad jqwik para \texttt{update} apoyada en un
  unico generador.
\end{itemize}

\section*{Conclusiones}
\addcontentsline{toc}{section}{Conclusiones}

\begin{itemize}
  \item \textbf{Randoop} es excelente como \emph{primera linea}: rapido, requiere casi nada
  de setup y suele exponer defectos obvios (campos sin inicializar, NPEs en operaciones
  publicas, contratos mal implementados). En cambio, sobre clases con dependencias externas
  fuertes o estado profundo, su exploracion aleatoria se queda corta.
  \item \textbf{EvoSuite} construye suites mas profundas porque su busqueda evolutiva
  optimiza directamente la cobertura. Paga el costo en complejidad de integracion (requiere
  un JDK compatible y su runtime puede ser fragil).
  \item \textbf{Mocking} con \texttt{EasyMock} es la herramienta natural cuando se quiere
  testear una clase \emph{aislada} de sus colaboradores: programar el doble es mucho mas
  practico que armar el entorno real.
  \item \textbf{Property-Based Testing con \texttt{jqwik}} permite expresar el comportamiento
  esperado en abstracto (``despues de \texttt{update()}, cada IP esta donde debe estar''), y
  el framework lo verifica sobre cientos de escenarios distintos.
  \item Las tres tecnicas son \emph{complementarias}, no excluyentes. Una buena suite las
  combina: PBT para invariantes, generacion automatica para descubrir casos, mocks para
  unidades aisladas.
\end{itemize}

\end{document}
